\documentclass[10pt,a4paper]{amsart}
\usepackage[margin=19mm]{geometry}
\usepackage{amsmath,amssymb,mathtools}
\usepackage[scr=rsfs,cal=euler]{mathalfa}
\usepackage{xfrac}
\usepackage[unicode,hidelinks,bookmarks=false]{hyperref}
\newcommand{\ie}{i.e.\ }
\newcommand{\cf}{cf.\ }
\newcommand{\resp}{respectively\ }
\newcommand{\Subsection}{\subsection}
\providecommand{\nobreakdash}{\nobreak\hbox{-}}
\setlength{\emergencystretch}{2em}
\multlinegap0pt
\usepackage{bm}
\usepackage{bbm}
\usepackage{dsfont}
\usepackage{xspace}
\usepackage{cancel}
\usepackage{mathtools}
\usepackage{amsthm}
\makeatletter
\@ifundefined{fed}{\newtheorem{fed}{Fed}}{}
\@ifundefined{lem}{\newtheorem{lem}[fed]{Lemma}}{}
\@ifundefined{pro}{\newtheorem{pro}[fed]{Proposition}}{}
\@ifundefined{teo}{\newtheorem{teo}[fed]{Theorem}}{}
\makeatother
\def\C{\mathbb{C}}
\def\R{\mathbb{R}}
\def\ele{\mathcal{L}}
\def\la{\lambda}
\def\w{\omega}
\newcommand{\supp}{{\rm supp\,}}
\title[Variations on two Cabrelli's works]{Variations on two Cabrelli's works}
\author{Elona Agora, Jorge Antezana, Diana Carbajal}
\date{Source: arXiv:2601.05422v1 (CC BY 4.0)}
\begin{document}
\maketitle

\noindent\emph{Source and licence.} Variations on two Cabrelli's works. Version arXiv:2601.05422v1 (CC BY 4.0), distributed under the Creative Commons Attribution 4.0 International licence, as recorded in the arXiv licence metadata for this exact version and shown on the abstract page https://arxiv.org/abs/2601.05422v1. Full rights notice: licenses/arxiv-2601.05422-CC-BY-4.0.txt.\par\smallskip

\noindent\textit{Editorial scope.} Counted unit: the Teo whose source label is \texttt{triangular} in the article (source0.tex lines 857--876, source label \texttt{triangular}), delivered together with the article's own complete original proof of it (source0.tex lines 878--929, one \texttt{proof} environment of 52 lines). No mathematics is written, repaired or re-derived: statement and proof are the article's own text line for line. Required context retained uncounted: range-properties (lines 741--751); An eigenvalue for R (lines 836--841). Every definition, display and label of those lines is retained unchanged. Omitted from this package and not redistributed: 1--740, 752--835, 842--856, 877--877, 930--1313. Nothing inside a retained excerpt is omitted or altered: every retained body is delivered exactly as the article's lines are written, so no omission receipt of any kind accompanies this package. Citations: the retained excerpts carry no citation command at all, so no bibliography entry of the article is transcribed and no reference is redistributed. Reference adaptation: none was needed and none was made; every reference inside the delivered unit resolves inside it, so no unresolved reference or citation is produced. Printed slips kept literal: no printed word, symbol, hypothesis or number of the counted statement or of its proof was altered. Layout and class: the article's own class is \texttt{amsart} with packages including amsfonts, amsthm, amsmath, latexsym, amssymb, mathtools, babel, inputenc, graphicx, upref, calligra, float, floatflt, xcolor; the wrapper is the lane's pinned \texttt{amsart} editorial wrapper at 10pt on A4, which restates the theorem-like environments the retained text needs and adds the article's own macro definitions that the retained text needs (\texttt{\textbackslash C}, \texttt{\textbackslash R}, \texttt{\textbackslash ele}, \texttt{\textbackslash la}, \texttt{\textbackslash supp}, \texttt{\textbackslash w}), with extra wrapper packages (bm, bbm, dsfont, xspace, cancel, mathtools). The delivered numbering is the unit's own. Assets: no retained span contains a figure environment, a graphics-inclusion command, a PDF-page inclusion, a table or a TikZ picture; no image file is copied into this package, the package declares no asset dependency and no image file is redistributed with it. Licence: version arXiv:2601.05422v1 of this article is distributed under the Creative Commons Attribution 4.0 International licence, as recorded in the arXiv licence metadata for this exact version and shown on the abstract page https://arxiv.org/abs/2601.05422v1. A full rights notice is delivered at licenses/arxiv-2601.05422-CC-BY-4.0.txt. Characters: every retained excerpt is pure ASCII, so the delivered engine, XeLaTeX with its default Latin Modern text font, reports no missing character.
	\begin{lem}\label{range-properties}
		Let $V, U$ be shift-invariant spaces of $L^2(\R^d)$ with associated range functions $J_V, J_U$, respectively. Then:
		\begin{enumerate}
			\item[i)]\label{item 1 range-properties} The orthogonal complement  $V^{\perp}$ is shift-invariant and $J_{V^{\perp}}(\w)=\left(J_V(\w)\right)^{\perp}$ for a.e. $\w\in I $.
			\item[ii)]\label{item 2 range-properties} The space  $V\cap U$ is  shift-invariant and its range function satisfies $J_{V\cap\, U}(\omega) = J_V(\omega) \cap J_U(\omega),$
			for a.e. $\omega\in I$.
			\item[iii)]\label{item 3 range-properties} If $V$ and $U$ are orthogonal, then $V\oplus U$ is shift-invariant and its range function is given by $J_{V\oplus U}(\omega) = J_V(\omega) \oplus J_U(\omega),$ for a.e. $\omega\in I$.
			\item[iv)] The length of $V$ can be characterized as follows:
			$$\mathcal{L}(V) = \underset{\omega\in I}{\text{\rm ess sup }}  \dim J(\omega).$$
		\end{enumerate}
	\end{lem}

	\begin{pro}\label{An eigenvalue for R}
		Let $V$ be a finitely generated shift-invariant space and let $L:V\rightarrow V$ be a shift-preserving operator with associated range operator $R$. Then, there exists a measurable function $\la: I\to \C$ such that $\la(\w)$ is an eigenvalue of $R(\w)$ for a.e. $\w\in I$. Moreover, there exists a function $\psi_\la \in V$ such that $\supp \| \mathcal T \psi_\la(\w)\| = \sigma(V)$ and
		\begin{equation}\label{eq:eigenfunction}
			R(\w)\mathcal T \psi_\la(\w)=\la(\w) \mathcal T \psi_\la(\w), \quad \text{for a.e.} \, \w \in I.
		\end{equation}	
	\end{pro}

	\begin{teo}\label{triangular}
		Let $L :V\rightarrow V$ be a bounded shift-preserving operator on a finitely generated shift-invariant space $V$ of length $\ell$.
		Then there exist shift-invariant subspaces of $V$, 
		$$
		V=V_\ell\supsetneq \dots\supsetneq V_{1},
		$$ 
		such that for every $1\leq j\leq\ell$,  $\ele(V_j)=j$, and each $V_j$ is \textit{invariant} for the operator $L$, that is, 
		$ L(V_j)\subseteq V_j$. Moreover, $V$ admits an orthogonal decomposition
		$$
		V = S(\psi_1)\oplus \dots \oplus S(\psi_\ell),
		$$
		where each $\psi_j$ is a Parseval frame generator of $S(\psi_j)$,  with the spectra satisfying the inclusions 
		$
		\sigma(S(\psi_{j+1}))\subset \sigma(S(\psi_j)),
		$
		and
		$$
		V_j = \bigoplus_{i=1}^{j} S(\psi_i), \quad 1\leq j \leq \ell.
		$$
	\end{teo}

	\begin{proof} We argue by induction on $\ell$.  If $\ell=1$, the result holds since $V$ is already a principal subspace.  
	Assume now that $\ell=n>1$, and that the result holds for all spaces of length strictly smaller than $n$. 
	Let $R$ be the range operator corresponding to $L$. By Proposition \ref{An eigenvalue for R}, 
	there exists a measurable function $\la_1 :I\to\C$ and a nonzero function $\psi_1\in V$ such that
	$$
	R(\omega)\mathcal T \psi_1(\omega) = \la_1(\omega)\,\mathcal T \psi_1(\omega), 
	\qquad \text{for a.e. }\omega \in I,
	$$
	and 
	$$
	\supp \|\mathcal T \psi_1(\cdot)\| = \sigma(V).
	$$
	In particular $\sigma(S(\psi_1))=\sigma(V)$. Thus $\psi_1$ defines the first generator of our decomposition. Moreover, without loss of generality, we may assume that
	$$
	\|\mathcal T\psi_1(\omega)\|=1, \qquad \text{for a.e. }\omega\in\sigma(S(\psi_1)),
	$$
	so that $\psi_1$ is a Parseval frame generator of $S(\psi_1)$.  
	The principal shift-invariant space
	$
	S(\psi_1)\subset V
	$
	is also invariant under $L$. Let $W:= V\ominus S(\psi_1)$. By Lemma \ref{range-properties}, $W$ is shift-invariant and the operator $L_W:W\to W$ defined by
	$$
	L_W := P_{W} L P_{W}
	$$
	is shift-preserving. Moreover, since $\text{span}\, \{ \mathcal{T} \psi_1 (\w) \}\neq \{0\}$ for a.e. $\w\in \sigma(V)$, and $\ele(V) = \ell$, we have that $\ele(W )= \ell-1$. 
	By the inductive hypothesis applied to $L_W:W\to W$, there exist orthonormal generators $\psi_2,\dots,\psi_\ell$ in $W$ such that
	$$
	W = S(\psi_2)\oplus\dots\oplus S(\psi_\ell),
	$$ 
	and the inclusions 
	$$
	\sigma(S(\psi_{i+1})) \subset \sigma(S(\psi_i)), \qquad 1\le i\le \ell-1,
	$$
	hold. Moreover, the spaces 
	$$
	W_j:= \bigoplus_{i=2}^j S(\psi_i), \quad  j\in \{2,\ldots,\ell\},
	$$
	are invariant under $L_W$ and, therefore invariant under $L$. Now define $V_1= S(\psi_1)$ and
	$$
	V_j= S(\psi_1) \oplus W_j=\bigoplus_{i=1}^{j} S(\psi_i), \quad  j\in \{2,\ldots,\ell\}.
	$$
	By construction, it follows that
	$$
	V = V_\ell \supsetneq V_{\ell-1} \supsetneq \cdots \supsetneq V_{1},
	$$
	where each $V_j$ is shift-invariant, invariant under $L$, and satisfies $\mathcal L(V_j)=j$. 
	Since  $\sigma(S(\psi_1))=\sigma(V)$,  $\sigma(S(\psi_2))\subseteq\sigma(S(\psi_1))$, we also have that
	$$
	\sigma(S(\psi_{i+1})) \subset \sigma(S(\psi_i)), \qquad 1\le i\le \ell-1.
	$$ 
\end{proof}

\end{document}
